\documentclass[12pt, a4paper]{article}

% --- PACKAGES---
\usepackage[utf8]{inputenc}       % Encodage (standard)
\usepackage[T1]{fontenc}          % Polices de sortie
\usepackage[english]{babel}        
\usepackage[backend=biber]{biblatex}
\usepackage{lmodern}              % Police plus lisible
\usepackage{microtype}            % Améliore l'apparence visuelle du texte
\usepackage{amsmath, amssymb, amsthm, amsfonts} 
\usepackage{mathtools}             % Extensions pour amsmath
\usepackage{mathrsfs}              % Lettres calligraphiques (\mathscr{L})
\usepackage{bm}                    % Symboles en gras (\bm{x})
\usepackage{empheq}
\usepackage{csquotes}
\usepackage{comment}

% --- CONFIGURATION DES THÉORÈMES ---
\newtheorem{theorem}{Theorem}
\newtheorem{lemma}[theorem]{Lemma}
\newtheorem{proposition}[theorem]{Proposition}
\newtheorem{corollary}[theorem]{Corollary}
\newtheorem{definition}[theorem]{Definition}
\newtheorem{example}[theorem]{Example}
\newtheorem{remark}[theorem]{Remark}

% --- MISE EN PAGE ---
\usepackage[margin=2.5cm]{geometry} % Marges
\usepackage{hyperref}               % Liens cliquables (sommaire, citations)
\hypersetup{colorlinks=true, linkcolor=blue, citecolor=red}

% --- TES MACROS (Raccourcis) ---
\newcommand{\R}{\mathbb{R}}
\newcommand{\N}{\mathbb{N}}
\newcommand{\C}{\mathbb{C}}
\newcommand{\E}{\mathbb{E}}

\newcommand{\diff}{\mathrm{d}}     % Pour un joli 'd' de différentielle
\newcommand{\dG}{\mathrm{d}\Gamma^*_x}
\newcommand{\dGb}{\mathrm{d}\Gamma^*}
\newcommand{\D}{\mathscr{D}}

\newcommand{\G}{\Gamma^*}
\newcommand{\mn}{\mathbf{n}}
\newcommand{\mm}[1]{\mathbf{m}}
\newcommand{\abs}[1]{\left| #1 \right|}
\newcommand{\norm}[1]{\left\| #1 \right\|}
\newcommand{\dd}{\,\mathrm{d}}
\newcommand{\z}[1]{\mathbf{z}_{#1}}

\newcommand{\gdb}[1]{\begingroup\color{red}GDB:\ #1\endgroup}
\newcommand{\lz}[1]{\begingroup\color{blue}LZ:\ #1\endgroup}


% --- INFOS DU MÉMOIRE ---
\title{Logical Order of the Proof}
\author{Gabriel de Boisgrollier}
\date{\today}

\addbibresource{mabibliozotero.bib}


\begin{document}
\maketitle

The goal is to prove, for all multi-indices $\beta$ that :

\begin{empheq}[left=\empheqlbrace]{align*}
    \E^{\frac{1}{p}} \abs{\Pi_{x\beta r}(x)}^p &\lesssim r^{\abs{\beta}}, \\
    \E^{\frac{1}{p}} \abs{\Pi_{x\beta r}^-(x)}^p &\lesssim r^{\abs{\beta}-2}, \\
    \E^{\frac{1}{p}} \abs{(\Gamma^*_x)^\gamma_\beta}^p &\lesssim x^{\abs{\beta}-\abs{\gamma}} \quad \text{for all populated }\gamma,
\end{empheq}
where the implicit constants only depend on $d$, $s$, $\beta$, $\gamma$, $p$, on $\psi$ only through the semi-norms, but not on $x$, $r$, and $\rho$.

To do, we proceed by induction on the multi-indices $\beta$, with the order $\abs{\cdot}_{\prec}$.
Furthermore, we need to enrich our induction hypothesis.
For $\beta$ a multi-index, let our set of hypothesis be :

\begin{itemize}
    \item Target bounds :
    \begin{empheq}[left=\empheqlbrace]{align}
    \E^{\frac{1}{p}} \abs{\Pi_{x\beta r}(x)}^p &\lesssim r^{\abs{\beta}}, \label{induction_plus}\\
    \E^{\frac{1}{p}} \abs{\Pi_{x\beta r}^-(x)}^p &\lesssim r^{\abs{\beta}-2}, \label{induction_minus}\\
    \E^{\frac{1}{p}} \abs{(\Gamma^*_x)^\gamma_\beta}^p &\lesssim x^{\abs{\beta}-\abs{\gamma}} \quad \text{for all populated }\gamma, \label{coeff}
    \end{empheq}
    \item Coherence bound : (the denomination will be clearer later)\begin{equation}\label{coherence}
\begin{aligned}
&\text{for all populated } \gamma,\quad
\left(
\int_{B_R} dx\,
\mathbb{E}^{\frac{2}{q}}
\left|
\left(
\mathrm{d}\Gamma_{x+y}^*
-\mathrm{d}\Gamma_x^*\Gamma_{x\,x+y}^*
\right)_\beta^\gamma
\right|^q
\right)^{\frac{1}{2}}
\\[6pt]
&\quad\lesssim
\begin{cases}
|y|^{\alpha-|\gamma|_p+\frac{D}{2}}(|y|+R)^{|\beta|-|\gamma|+|\gamma|_p-\alpha},& \text{if } \alpha-|\gamma|_p+\frac{D}{2}>0,
\\[6pt]
(|y|+R)^{|\beta|-|\gamma|+\frac{D}{2}},
& \text{else}.
\end{cases}
\end{aligned}
\end{equation}
\item Estimate on coefficients:
\begin{equation}\label{dcoeff_estimate}
    \text{for all populated }\gamma,\quad\left(\int_{B_R} \dd x \mathbb{E}^{\frac{2}{q}}\abs{(\dG)^\gamma_\beta}^q \right)^\frac{1}{2} \lesssim R^{\frac{D}{2}+\abs{\beta} - \abs{\gamma}}, \\
\end{equation}
\end{itemize}

The proof relies on the spectral gap inequality :
\begin{equation}
        \E^{\frac{1}{p}} \abs{\Pi^-_{\beta r}(0)}^p \lesssim \abs{\E \Pi^-_{\beta r}(0)} 
        + \sup \{\E^{\frac{1}{q}} \abs{\delta \Pi^-_{\beta r}(0)}^q ; \delta \xi \text{ random with }\E \norm{\delta \xi}^{p^*}_{\dot{H}_s} \leq 1\},
\end{equation}
for $1\leq q < p^* \leq 2 \leq p$.
We will prove estimates for all $p$, $q$, $p^*$ in this range so we will be able to use the induction hypothesis with $p'\geq p$ and $p'^*>q'\geq q$ as it naturally occurs with Hölder's inequality.

We claim that we can control $\abs{\E \Pi^-_{\beta r}(0)}$  when $\abs{\beta}<2$.
We need to control the second term.




Sketch of the reasoning :

A promising strategy is to express the increments as follows:
\begin{align*}
    \Pi^-_{x+y\beta} - \Pi^-_{x\beta} &= ((\operatorname*{Id} - \Gamma_{x x+y}^*)\Pi^-_{x+y})_\beta \\
    &= \sum_{\gamma} (\operatorname*{Id} - \Gamma_{x x+y}^*)^\gamma_\beta \Pi^-_{x+y\gamma} \\
    &= - \sum_{\gamma \prec \beta} (\Gamma_{x x+y}^*)^\gamma_\beta \Pi^-_{x+y\gamma} \\
    &= - \sum_{\substack{\gamma \prec \beta \\ \gamma \neq \text{pp}}} (\Gamma_{x x+y}^*)^\gamma_\beta \Pi^-_{x+y\gamma},
\end{align*}
taking advantage of the triangularity of the operator $\operatorname*{Id} - \Gamma_{x x+y}^*$ with respect to $\prec$, and the fact that $\gamma = \text{pp}\implies \Pi^-_{z\gamma} = 0$.
As a result, we can use the induction hypothesis \eqref{induction_minus} to estimate $\Pi^-_{x+y\gamma}$.

In addition, for all $z\in \R^{d}$ and $\gamma\neq \text{pp}$,
\begin{equation}
    (\Gamma_{z}^*)_\beta^\gamma = \sum_{k\delta_3 + \beta_1 + \dots + \beta_l = \beta} (\Gamma^*_z)_{\beta_1}^{\delta_{\mn_1}}\cdots (\Gamma^*_z)_{\beta_l}^{\delta_{\mn_l}},
\end{equation}
and $k\delta_3 + \beta_1 + \dots + \beta_l = \beta \implies \forall i\in \{1,\dots,l\}, \, \beta_i \prec \beta$,
so that \eqref{coeff} is true for $\gamma\neq \text{pp}$ : 
\begin{align*}
    \E^{\frac{1}{p}} \abs{(\Gamma^*_z)^\gamma_\beta}^p &\lesssim \sum \E^{\frac{1}{lp}} \abs{(\Gamma^*_z)_{\beta_1}^{\delta_{\mn_1}}}^{pl} \cdots \E^{\frac{1}{lp}} \abs{(\Gamma^*_z)_{\beta_l}^{\delta_{\mn_l}}}^{pl} \\
    &\lesssim z^{\abs{\beta_1}+\cdots+\abs{\beta_1} - \abs{\mn_1}-\cdots - \abs{\mn_l}}\\
    &\lesssim z^{\abs{\beta}-\abs{\gamma}}.
\end{align*}







Fix $\beta \neq 0$ such that $\forall \gamma \prec \beta,$ $(P_\gamma)$ holds.
Fix $p$, $q$ such that $1\leq q < p^* \leq 2 \leq p$.



From \eqref{coherence}, we deduce that the germ $(\delta \Pi_\beta^- - (\dG \Pi_x^-)_\beta)$ is coherent. By reconstruction, we get :
\begin{equation}
        (\int_{B_R} \dd x \mathbb{E}^{\frac{2}{q}}\abs{(\delta \Pi^- - \dG \Pi_x^-)_{\beta r}(x)}^q)^\frac{1}{2} \lesssim r^{\alpha - 2 + \frac{D}{2}}(r+R)^{\abs{\beta} - \abs{\alpha}}. \label{reconstruction}
\end{equation}
By \eqref{dcoeff_estimate} and \eqref{induction_minus}, the germ $(\dG \Pi_x^-)_\beta$ is homogeneous, so we get by triangle inequality an estimate on $\delta \Pi^-$:
\begin{equation*}
        (\int_{B_R} \dd x \mathbb{E}^{\frac{2}{q}}\abs{\delta \Pi^-_{\beta r}(x)}^q)^\frac{1}{2} \lesssim r^{\alpha - 2 + \frac{D}{2}}(r+R)^{\abs{\beta} - \abs{\alpha}}.
\end{equation*}
By Sobolev's inequality, we can post-process this inequality to obtain the pointwise estimate~:
\begin{equation}
    \label{control_deltaPi^-}
    \E^{\frac{1}{q}} \abs{\delta \Pi^-_{\beta r}(x)}^q \lesssim r^{\alpha - 2 - \frac{D}{2}}(r+\abs{x})^{\abs{\beta}-\alpha + \frac{D}{2}}.
\end{equation}
We have the spectral gap inequality :
\begin{equation}
        \E^{\frac{1}{p}} \abs{\Pi^-_{\beta r}(0)}^p \lesssim \abs{\E \Pi^-_{\beta r}(0)} 
        + \sup \{\E^{\frac{1}{q}} \abs{\delta \Pi^-_{\beta r}(0)}^q ; \delta \xi \text{ random with }\E \norm{\delta \xi}^{p^*}_{\dot{H}_s} \leq 1\},
\end{equation}
and for $\abs{\beta}<2$, we have the control :
$$
\abs{\E \Pi^-_{\beta r}(0)}  \lesssim r^{\abs{\beta}-2}.
$$
Estimate \eqref{control_deltaPi^-} at $x=0$ with these last two estimates gives :
\begin{equation}
    \E^{\frac{1}{p}} \abs{\Pi_{\beta r}^-(0)}^p \lesssim r^{\abs{\beta}-2}, \\
\end{equation}
which generalizes to the target estimate \eqref{induction_minus} because $\Pi_x^- (\cdot + x)\overset{\mathrm{law}}{=}\Pi^-$.
For $\abs{\beta}>2$,
we need estimate \eqref{coeff} for $\gamma \neq $pp, which we prove by an algebraic argument





By integration, we deduce target estimate \eqref{induction_plus}













\end{document}